\scanpage{62}
Comme $\ell_\bullet^G=0$, le triangle exact de dualité pour $G$
se réduit à
\diagram{\begin{tikzcd}[ampersand replacement=\&,column sep=small,row sep=small]
\&\check\ell_\bullet^{G^*}[1]\arrow[dl]\&\\
0\arrow[rr,no head]\&\&\Delta_\bullet(G)\arrow[ul,"\unsure{\sim|2}"']
\end{tikzcd}}
i.e. ne donne pas de structure supplémentaire
\deleted{\unsure{rel. à} $\unsure{G}$} \inserted{en sus}
du cristal $M$. On a
\[
\begin{aligned}
\check\ell_\bullet^{G^*}[1]
&\simeq\unsure{\underline{\Sigma}^*(G)|\underline{\mathbb D}^*(G)}
\simeq\varepsilon^*(M_{/p})
=\deleted{\unsure{\varepsilon}}
R\underline{\operatorname{Hom}}_{\mathbb Z}(G,\underline{\mathcal O}_S)\\
&\simeq\check G\overset L\otimes_{\mathbb Z}\underline{\mathcal O}_S
\simeq(\check G\overset L\otimes_{\mathbb Z}\mathbb F_p)
\otimes_{\mathbb F_p}\underline{\mathcal O}_S.
\end{aligned}
\]

\uline{Groupes de type multiplicatif}
(\unsure{annulés} par puiss. $p^n$ de $p$)

Soit $G$ \deleted{est} un tel, il est défini par $G^*$,
groupe étale annulé par $p^n$, auquel s'appliquent les réflexions
précédentes. On aura
\[
\mathbb D^*(G)=\mathbb D_*(G^*)
\]
i.e.
\[
\mathbb D^*(G)_{S'}=(G^*_{S'})\otimes_{\mathbb Z_p}
\underline{\mathcal O}_{S'}
\]
\inserted{$G^*_{S'}$: groupe étale sur $S'$ prolongeant $G^*$}

Ici
\[
V:\mathbb D^*(G)\longrightarrow\mathbb D^*(G^{(p)})
=\mathbb D^*(G)^{(p)}
\]
est un isomorphisme, savoir
$\mathbb D_*(f_{G^*}:G^*\longrightarrow G^{*(p)})$,
et $F$ est \inserted{par suite} déterminé uniquement en termes de $V$.

On a encore des équivalences de catégories
\diagram{\[\begin{array}{ccl}
\text{groupes de B.T. tronqués d'échelon }n\text{ de t.m.}
&\simeq&
\text{cristaux loc. libres de t.f.}\\
&&\text{pour }S/\Lambda_n,\text{ avec }V\\
&&\text{un isom (et }F\text{ si on veut)}
\end{array}\]}
\margin{72}
